\documentclass[12pt,a4paper]{article}

% ============================================================
% MediBridge Project Proposal
% Recreated in LaTeX from the supplied PDF
% ============================================================

\usepackage[a4paper,margin=1in]{geometry}
\usepackage{mathptmx}
\usepackage{setspace}
\usepackage{enumitem}
\usepackage{titlesec}
\usepackage{parskip}
\usepackage[hidelinks]{hyperref}

% Match the compact report appearance
\setstretch{1.08}
\setlength{\parindent}{0pt}
\setlength{\parskip}{0.65em}

\setlist[itemize]{
    leftmargin=1.5em,
    itemsep=0.35em,
    topsep=0.25em
}

\setlist[enumerate]{
    leftmargin=1.8em,
    itemsep=0.35em,
    topsep=0.25em
}

% Section numbering/spacing
\titleformat{\section}
  {\normalfont\Large\bfseries}
  {\thesection}
  {0.65em}
  {}

\titleformat{\subsection}
  {\normalfont\large\bfseries}
  {\thesubsection}
  {0.55em}
  {}

\titlespacing*{\section}{0pt}{1.0em}{0.45em}
\titlespacing*{\subsection}{0pt}{0.7em}{0.3em}

% ============================================================

\begin{document}

% ------------------------- TITLE PAGE -------------------------

\begin{center}

\vspace*{0.15cm}

{\LARGE\bfseries MediBridge: Healthcare Access and Financial Assistance Platform}

\vspace{1.15cm}

{\large\bfseries KUNAL MOTWANI}

\vspace{0.08cm}

{\large 1024160137}

\vspace{0.35cm}

{\large\bfseries SURYANSH PAHUJA}

\vspace{0.08cm}

{\large 1024160058}

\vspace{0.65cm}

{\large August 21, 2026}

\end{center}

\vspace{0.7cm}

% ------------------------- CONTENT ----------------------------

\section{Higher-order Goal}

Healthcare accessibility remains a major challenge due to fragmented
information, expensive treatment options, and difficulty in arranging
financial support. This project aims to build MediBridge, a software
platform that connects healthcare discovery and financial assistance into
a unified ecosystem.

The engineering objective is to design a scalable web platform where
patients can discover hospitals, manage medical information, understand
treatment options, and access financial support workflows. The system
focuses on improving decision-making and reducing delays during critical
healthcare situations.

\section{Time-to-value}

The project follows an MVP-first development strategy:

\begin{itemize}
    \item First iteration delivers patient registration, healthcare
    profiles, and hospital discovery.
    \item Subsequent iterations introduce medical record management and
    financial assistance workflows.
    \item Continuous testing and user feedback guide feature improvements.
\end{itemize}

\section{Problem Statement}

Patients often face difficulty finding affordable healthcare services.
Information about hospitals, treatment costs, doctor availability, and
financial assistance is usually distributed across multiple platforms.

Existing healthcare systems frequently focus on individual services
rather than providing a complete patient journey. Patients may struggle
with maintaining medical records, comparing options, and arranging funds
for treatment.

This project addresses the gap by developing a centralized healthcare
assistance platform combining medical information access and financial
support mechanisms.

\section{Proposed Solution}

\subsection{Overview}

MediBridge provides:

\begin{itemize}
    \item Patient dashboard.
    \item Hospital and doctor search.
    \item Medical record storage.
    \item Treatment cost comparison.
    \item Loan and fundraising assistance workflows.
    \item Emergency health profile.
\end{itemize}

\subsection{Core Workflow}

\begin{enumerate}
    \item Patient creates a secure profile.
    \item User searches treatment options and hospitals.
    \item Medical records are uploaded and managed.
    \item Financial assistance requests can be initiated.
    \item Emergency information can be accessed when required.
\end{enumerate}

\section{Solution Approach}

\subsection{System Architecture}

The platform follows a three-tier architecture:

\begin{itemize}
    \item \textbf{Frontend:} Responsive patient and administrator
    interfaces.
    \item \textbf{Backend:} REST APIs for authentication, workflows,
    and business logic.
    \item \textbf{Database:} Stores users, hospitals, records, and
    application history.
\end{itemize}

\subsection{Security Approach}

The system will implement authentication, authorization, secure
document handling, and privacy-aware data storage.

\subsection{CI/CD and Testing}

Automated testing and deployment pipelines will enable reliable
incremental development.

\section{Evaluation Criterion}

The system will be evaluated using:

\begin{itemize}
    \item Patient workflow completion time.
    \item Recommendation usefulness.
    \item Document access reliability.
    \item User satisfaction scores.
    \item System availability.
\end{itemize}

\section{Scalability}

The architecture supports future growth through modular services, cloud
deployment, database indexing, and API-based integrations with hospitals
and financial organizations.

\section{Engine Availability Heuristic}

The project can be implemented using mature web frameworks, cloud
databases, authentication systems, storage services, testing tools, and
deployment platforms.

\section{Project Scope and Deliverables}

\subsection{Initial Deliverable}

\begin{itemize}
    \item Authentication system.
    \item Patient dashboard.
    \item Hospital search.
    \item Medical profile management.
\end{itemize}

\subsection{Subsequent Deliverables}

\begin{itemize}
    \item Financial assistance workflow.
    \item Recommendation engine.
    \item Notification system.
    \item Partner integrations.
\end{itemize}

\section{Risks and Mitigations}

Major risks include healthcare data privacy, large scope, and dependency
on external healthcare information. These will be handled through
limited MVP scope, secure storage practices, and modular development.

\section{Summary}

MediBridge proposes a scalable healthcare support platform that combines
healthcare discovery and financial assistance into a single
patient-focused solution.

\end{document}
